\uuid{bNTK}
\titre{Équation droite/plan}
\auteur{bodin ; hanani ; mzari}
\organisation{exo7}
\licence{CC BY-NC-SA 4.0}
\datecreate{2020-01-20}
\niveau{L1}
\module{Géométrie}
\chapitre{Géométrie de l'espace}
\sousChapitre{Plans}
\difficulte{2}
\theme{Plans -- Droites, équation droite, plan}
\exo7qcmid{lille-1.geometrie-espace.010}

\contenu{
\texte{
Pour ces questions, l'espace est muni d'un repère orthonormé direct  $(O, \vec{i}, \vec{j}, \vec{k})$.
}

\begin{question}
\qclassification{141.03,141.04}
\qtags{motcle=équation droite/plan}
Soit $P : x+y-z=0$,  $P' : x-y=2$ et $P'' : y-z=3$ trois plans. L'intersection de ces trois plans est :
\begin{answers}

    \bad{Vide.}

    \bad{Une droite.}

    \good{Un point.}

    \good{Le point de coordonnées $(-3,-5,-8)$.}

\end{answers}
\begin{explanations}
On résout le système constitué des équations des trois plans.
\end{explanations}
\end{question}
}
